\chapter{绪论}

\section{任务背景}

随着大语言模型和多模态大模型的发展，面向文档的智能问答系统已经从简单的文本检索逐步走向版面理解、图表推理和可追溯证据定位。大量真实文档并不是线性文本，而是由标题、段落、表格、图像、公式、脚注、页面布局和视觉标记共同组成。对于 PDF、扫描文档、课程讲义、财报、论文和作业说明等材料，关键信息经常以表格行列、图表趋势、局部截图或跨页上下文的形式出现。如果仅将文档抽取为纯文本，再使用传统 RAG 检索，很容易出现三类问题：一是 OCR 或文本抽取丢失视觉布局，二是图表数值和视觉元素无法可靠转写，三是生成答案缺乏可检查的页面证据。

本课程大作业给出的题目 3 要求构建“文档理解+多模态检索问答系统”：系统应支持上传多页 PDF 或图像，自动分块文本、表格与图形；支持自然语言问题查询，并在回答中标出引用页面或图号；在 DocVQA、ChartQA 子集或自建数据上评测；并对比纯文本 RAG 与多模态 RAG 的效果差异。期末报告还要求提交最终报告、系统实现细节、代码仓库、定量实验结果、可视化/错误分析以及反思与展望。本文即围绕这些要求，对本项目的系统设计、工程实现、实验评测和局限进行完整总结。

\section{问题定义}

给定一个多页文档 $D=\{p_1,p_2,\ldots,p_n\}$ 与用户问题 $q$，系统需要从文档中检索出支持回答的证据集合 $E=\{e_1,e_2,\ldots,e_k\}$，并生成自然语言答案 $a$。与普通文本问答不同，本文关注的证据既包括文本片段，也包括页面图像、局部视觉区域、表格行列和图表区域。因此，一个证据单元不仅包含文本内容，还应包含页码、chunk 标识、区域坐标、图像路径和来源类型。系统输出应满足以下目标：

\begin{itemize}
  \item \textbf{正确性}：回答应尽量与标准短答案或人工期望答案一致。
  \item \textbf{可追溯性}：每个答案应标明对应页码、chunk 或视觉区域，便于人工检查。
  \item \textbf{多模态性}：当问题涉及图表、图片、表格或视觉布局时，应使用视觉 evidence，而不是仅依赖 OCR 文本。
  \item \textbf{交互性}：用户可以上传文档、查看解析后的知识库 chunk，并在 PDF 预览中看到对应高亮区域。
\end{itemize}

\section{开题目标与完成概览}

开题报告中，本项目计划实现三个核心目标：第一，构建多模态 RAG 架构，将文本、表格、图像等信息组织为可检索证据；第二，设计路由决策机制，根据问题类型选择合适的多模态模型或处理路径；第三，提供可解释性溯源，在回答中标注页码、图表或视觉区域。期末阶段，项目围绕这三个目标完成了一个可运行的系统原型，并对部分原计划进行了工程化调整。

总体上，文档解析、多模态 evidence、检索重排、前端上传、PDF 预览、高亮溯源和 DocVQA/ChartQA 评测均已完成；系统还补充了 text、table、figure、code 等 chunk 类型分离、表格 Markdown 入库、统一文本-图像 embedding 接口和 GRPO-style auto router。受课程周期限制，GRPO 路由没有进行真实强化学习训练，而是提供候选模型打分、策略更新钩子和可运行的自动路由入口。这样的调整使系统能够在课程周期内形成稳定闭环，并保留后续继续扩展学习式路由的接口。

\section{系统贡献}

本项目最终实现了一个可运行的多模态文档问答原型，代码仓库包含后端 RAG pipeline、评测脚本、FastAPI 服务层、React 前端和 Gradio 兼容入口。相比开题阶段的系统设想，最终版本从“方案设计”推进到“可评测、可交互、可解释”的工程闭环。本文主要贡献如下：

\begin{enumerate}
  \item \textbf{实现端到端文档 RAG 流程}：系统支持 PDF/图片输入，完成页面渲染、结构化 chunk 构建、向量索引、重排、路由和生成式回答。
  \item \textbf{引入多模态 evidence}：在文本 evidence 之外，系统加入表格 Markdown、页面级视觉摘要、图像区域和局部高清 crop，并将视觉区域作为可检索证据参与 MM-RAG。
  \item \textbf{构建真实模型栈}：后端支持 BGE-M3 embedding、FAISS 索引、BGE-reranker-v2-m3 重排和 Qwen3-VL-8B-Instruct 推理，同时保留 mock backend 便于无模型环境下开发调试。
  \item \textbf{实现可解释前端交互}：React 前端支持上传、解析轮询、知识库 chunk 列表、连续 PDF 预览、chunk 选择与页面高亮联动，接近 RAGFlow 类文档知识库的使用体验。
  \item \textbf{完成定量评测与错误分析}：在 20 个 DocVQA 与 20 个 ChartQA 样本上比较 Text-RAG 与 MM-RAG，MM-RAG 总体 EM 达到 0.9750，明显优于 Text-RAG 的 0.3750。
\end{enumerate}

\section{报告结构}

本文后续结构如下：第二章按照 Embedding/Fusion、向量检索、Reranker、GRPO 路由、VLM 与 PDF 高亮等组件说明相关技术与功能对应；第三章给出系统方法与工程实现；第四章介绍实验设置、数据集、指标和运行环境；第五章展示实验结果、案例与错误分析；第六章总结项目收获，并从模型局限和 AGI 视角讨论未来方向。
